\begin{examplebox}
\textbf{例1：先验证点在曲线上，再验证乘积。}
椭圆 $x^2/4+y^2/3=1$，取 $P=(1,3/2)$。
代入得 $1/4+(9/4)/3=1$，点合法。焦点 $F_1=(-1,0)$、$F_2=(1,0)$，
于是 $r_1=5/2,r_2=3/2$。切线为 $x+2y=4$，原点到切线距离 $d=4/\sqrt5$。
\[
\sqrt{r_1r_2}\,d=\frac{\sqrt{15}}2\frac4{\sqrt5}=2\sqrt3=ab.
\]
原坐标 $(1,\sqrt3/2)$ 代入椭圆仅得 $1/2$，不能作为切点。
\par\medskip\noindent
\textbf{例2：距离反求焦半径乘积。}
若 $a=5,b=3,d=4$，则 $r_1r_2=(ab/d)^2=225/16$。
另由 $r_1+r_2=10$ 得 $\{r_1,r_2\}=\{5-5\sqrt7/4,5+5\sqrt7/4\}$。
因为各焦半径须在 $[1,9]$ 内，上式中较小值约 $1.693$，满足范围，条件可实现。
\end{examplebox}
